\documentclass[11pt]{amsart}
\usepackage{amssymb,mathtools}
\usepackage[margin=1.15in]{geometry}
\usepackage[hidelinks]{hyperref}
\newtheorem{theorem}{Theorem}[section]
\newtheorem{lemma}[theorem]{Lemma}
\theoremstyle{remark}\newtheorem{remark}[theorem]{Remark}
\newcommand{\Hom}{\operatorname{Hom}}
\newcommand{\Spec}{\operatorname{Spec}}
\newcommand{\Nat}{\operatorname{Nat}}
\newcommand{\Alg}{\operatorname{Alg}}
\newcommand{\UU}{\mathcal U}
\title[Full faithfulness of the functor of points]{Full faithfulness of the functor of points\\An informal proof and its interfaces}
\date{October 29, 2026}
\begin{document}\maketitle

\section{Statement and size conventions}
Fix a commutative ring $k$ with carrier in $\UU_\ell$. Test algebras are
\emph{all} commutative $k$-algebras with carriers in $\UU_\ell$.
Their object type belongs to $\UU_{\ell+1}$ and their hom sets to
$\UU_\ell$. We allow functors and natural transformations on this test
category to live in the larger universe. No smallness theorem for all
these algebras is asserted or needed.

Schemes are qcqs locally ringed lattices with lattices and section rings
at level $\ell$. A $k$-structure is a local morphism to $\Spec k$;
morphisms over $k$ commute with these maps. For such a scheme $X$ put
\[
 h_X(B)=\Hom_k(\Spec B,X),\qquad
 h_X(h)(p)=p\circ\Spec h\quad(h:B\to B').
\]
These are sets. The identity and composition laws for $\Spec$ make
$h_X$ a covariant set-valued functor.

\begin{theorem}\label{thm:main}
For qcqs $k$-schemes $X,Y$, postcomposition is an equivalence
\[
 \Hom_k(X,Y)\longrightarrow\Nat(h_X,h_Y).
\]
\end{theorem}
The proof below constructs its inverse. It does not characterise which
functors are represented by schemes. In fact only the source $X$ needs
the finite affine-cover hypothesis; the descent argument works for any
locally ringed lattice $Y$ of the stated size.
The construction first establishes the absolute statement on all small
ring tests. Section~5 passes from that statement to tests over $k$.

\section{Finite-cover descent for morphisms}
Let $X$ have lattice $M$, and let $e_1,\ldots,e_n$ have join $\top$.
For each $i$ let $f_i:X|e_i\to Y$ be a local morphism. For a compact
open $v$ of $Y$, write $p_i(v)$ for its inverse image, viewed as an
element of $M$ below $e_i$. Agreement on overlaps means
\begin{equation}\label{eq:opens}
 p_i(v)\wedge e_j=p_j(v)\wedge e_i,
\end{equation}
and agreement of the corresponding section maps after restriction to
these intersections. This is equality of the restricted morphisms,
expressed without identifying different presentations of a slice.

\begin{lemma}[Separation]\label{lem:separation}
Two morphisms $X\to Y$ agreeing on every $X|e_i$ are equal.
\end{lemma}
\begin{proof}
In a distributive lattice, $a=\bigvee_i(a\wedge e_i)$, so their inverse
images of each $v$ agree. Once this equality is known, their two images
of a section over $v$ lie in the same section ring. Their restrictions
to the finite cover by the intersections with the $e_i$ agree; separation
of the structure sheaf gives equality. Extensionality of ringed-lattice
maps, and the fact that locality is a proposition, give the result.
\end{proof}

\begin{lemma}[Gluing]\label{lem:gluing}
Compatible $f_i$ have a unique local amalgam $f:X\to Y$.
\end{lemma}
\begin{proof}
Define the inverse-image function by
\[
 p(v)=\bigvee_i p_i(v).
\]
It preserves bottom and joins, and sends top to $\bigvee_i e_i=\top$.
Equation~\eqref{eq:opens} and $p_i(v)\leq e_i$ give
\[
 p(v)\wedge e_i=p_i(v).
\]
Indeed each $p_j(v)\wedge e_i=p_i(v)\wedge e_j$ is below $p_i(v)$,
and the $i$th term gives the converse inequality. Testing on the cover
$e_i$ now proves preservation of meets, since each $p_i$ preserves meets.

For $s\in\mathcal O_Y(v)$, the sections $f_i^\sharp(s)$ over $p_i(v)$
agree on pairwise intersections. Glue them to a section over $p(v)$.
This defines $f^\sharp_v$. Preservation of ring operations follows by
separation on the same cover, because each $f_i^\sharp$ is a ring map.
Naturality follows by restriction and separation on the cover of $p(w)$
when $w\leq v$.

For locality, restriction of invertibility opens gives
\[
 D_{p(v)}(f^\sharp s)\wedge e_i
   =D_{p_i(v)}(f_i^\sharp s)=p_i(D_v(s)).
\]
The first equality uses $D_{p(v)}(f^\sharp s)\leq p(v)$, and the
second is locality of $f_i$. These are precisely the intersections of
$p(D_v(s))$ with the $e_i$. Lattice separation gives locality of $f$.
The constructed morphism restricts to every $f_i$ by its defining
equations; uniqueness is Lemma~\ref{lem:separation}.
\end{proof}

Restriction to an open and to a smaller open must also be identified
canonically: $(X|e)|u\cong X|u$ for $u\leq e$. For a morphism $g:Z\to X$,
restriction to the inverse image of $e$ gives $Z|g^*e\to X|e$.
These facts follow directly from the slice universal property. All
equalities of section maps can be checked after restriction in $X$;
no transport of a whole locally ringed lattice is necessary.

\section{Affine refinements}
\begin{lemma}\label{lem:basis}
Every compact open of a qcqs scheme has a merely existing finite cover
by affine opens. Every inverse image of a finite cover is a finite cover.
\end{lemma}
\begin{proof}
Choose a finite affine cover $e_i$ only while proving this mere existence.
For a given compact open $u$, the opens $u\wedge e_i$ cover it. Through
the affine chart $X|e_i\cong\Spec A_i$, each corresponds to an element of
the Zariski lattice of $A_i$. Such an element merely has a finite
presentation as a join of basic opens. Finite choice under truncation
combines these presentations. Basic opens are affine by
$(\Spec A_i)|D(a)\cong\Spec((A_i)_a)$, and the slice identifications above
make these affine opens of $X$. Flattening the finite families gives
the desired cover. The second assertion follows because inverse-image
lattice maps preserve finite joins and top.
\end{proof}

Consequently a pair of morphisms from an open agrees whenever it agrees
after every affine probe into that open: choose an affine cover and
apply Lemma~\ref{lem:separation}. The use of a mere cover is legitimate
since equality of morphisms is a proposition.

\section{Constructing a morphism from a natural transformation}
Let $\alpha:h_X\to h_Y$ be natural. First note the following affine
calculation. For an affine source $\Spec A$ and a natural transformation
$h_{\Spec A}\to h_Y$, evaluation at $1_{\Spec A}$ produces a morphism
$\Spec A\to Y$. Every point $q:\Spec B\to\Spec A$ over $k$ is
$\Spec h$ for a unique algebra map $h:A\to B$, by the universal property
of $\Spec$ and $\Gamma(\Spec B)\cong B$. Naturality therefore determines
the transformation on $q$ by precomposition. This argument uses no
small object type for the category of all test algebras.

Choose affine charts $j_i:U_i\to X$, with identifications
$U_i\cong\Spec A_i$, and put
\[
 f_i=\alpha_{A_i}(j_i):U_i\to Y.
\]
For every affine probe $q:\Spec B\to U_i$, naturality gives
\begin{equation}\label{eq:natural}
 f_i\circ q=\alpha_B(j_i\circ q).
\end{equation}
For a probe into $U_i\cap U_j$, the right side is the same for $i$ and
$j$. The overlap has a finite affine cover by Lemma~\ref{lem:basis};
Lemma~\ref{lem:separation} proves that $f_i$ and $f_j$ agree there.
Lemma~\ref{lem:gluing} gives $f:X\to Y$ over $k$. To check the
$k$-condition, compare its composite with $Y\to\Spec k$ to the structure
map of $X$ on the chart cover, and use separation.

For any $p:\Spec B\to X$, the inverse images of the $U_i$ cover
$\Spec B$. Refine these by finitely many basic affine opens. On such
an open, with inclusion $q:\Spec C\to\Spec B$, the composite $pq$
factors through some $j_i$. Equation~\eqref{eq:natural} and naturality
of $\alpha$ at $B\to C$ imply
\[
 f\circ p\circ q=\alpha_B(p)\circ q.
\]
Separation on this cover gives $fp=\alpha_B(p)$. Thus the transformation
induced by $f$ is $\alpha$. Conversely, if $\alpha$ is postcomposition
with a given $g$, its chart values are $gj_i$, whose unique amalgam is
$g$. These are the two inverse laws in Theorem~\ref{thm:main}.

The apparent choice of chart cover is harmless. The type
\[
 \sum_{f:\Hom_k(X,Y)}\ \prod_B\prod_{p:h_X(B)}(fp=\alpha_B(p))
\]
is a proposition: equality of its first coordinates is detected on an
affine cover, and the remaining coordinates are propositions. A cover
merely exists, so it may be eliminated into this type to obtain the
unique $f$. This also accounts for its larger universe.

\section{Passing from rings to algebras over the base}
Here is a precise reduction of the relative statement to the absolute
one. Let $s_X:X\to\Spec k$ and $s_Y:Y\to\Spec k$ be the structure maps.
For a test ring $B$, an ordinary point $p:\Spec B\to X$ determines a
unique ring map $\iota:k\to B$ whose spectrum is $s_Xp$. Equivalently,
forgetting the algebra structure is an equivalence
\[
 \sum_{\iota:\Hom(k,B)}h_X^k(B,\iota)\simeq h_X(B).
\]
For fixed $p$, the type of lifts on the left is a proposition, since
the affine correspondence makes $\iota$ unique and the equation for
being over $k$ is a proposition. It is inhabited by that correspondence.
Under a ring map $h:B\to C$, such a lift reindexes to the structure
$h\iota$; uniqueness identifies it with the canonical lift of the
reindexed ordinary point.

Given a transformation on $k$-algebras, apply it to these lifts and forget
the base equations. The uniqueness just proved and relative naturality
give a natural transformation on all rings. The absolute theorem recovers
$f:X\to Y$. On every affine probe $p$, the base equation of the output
says $s_Yfp=s_Xp$. Affine probes detect equality, so $s_Yf=s_X$.
Forgetting a relative point is injective, and every relative point is
the canonical lift of its underlying ordinary point. Hence this recovered
map induces the original relative transformation. The converse inverse
law follows from faithfulness, again using canonical lifts of ordinary
points. This proves the relative equivalence without replacing all
$k$-algebras by finitely presented ones.

\section{Proof review and implementation order}
The proof requires all test algebras to be small, not finitely presented:
chart rings, localisation rings, and rings of affine refinements are all
allowed. Each is still at level $\ell$. No chosen presentation, choice
of prime ideals, or resizing principle occurs. The use of a chart cover
is always under a proposition when independence matters.

The interfaces used in the proof are concrete: finite-cover separation and
gluing for morphisms; nested slices and restriction of morphisms to
inverse-image opens; transport of affine opens through chart
isomorphisms; and the finite affine-refinement lemma. The algebraic
identity and composition laws for $\Spec$ and a large-test-category
functor-of-points interface then express the final inverse laws.
The finite-presentation category and its Yoneda theorem are not used to
bypass any of these geometric steps.

\section{Connection with the implementation}
The supporting geometric lemmas are now checked in
\texttt{sag/qcqs/}. Finite-cover detection of opens and gluing of lattice
maps are in \texttt{LatticeCover} and \texttt{LatticeDescent}.
\texttt{MapAgreement} and \texttt{MapSeparation} prove separation of
morphisms, with the latter phrased using actual source restrictions.
\texttt{SheafMaps} glues ring maps on finite covers;
\texttt{MapGluing} constructs ringed morphisms, proves their restriction
equations and uniqueness, and shows that local inputs give a local
amalgam. Both the ringed and local amalgam types are contractible.

\texttt{SliceMap}, \texttt{SliceNested}, and \texttt{SliceIso} provide
restriction to inverse-image opens, the nested-slice isomorphism, and
restriction of isomorphisms. \texttt{AffineTransport} transfers
affineness. The affine case of refinement is proved in
\texttt{AffineBasis}. Finite families are concatenated in
\texttt{FiniteLatticeFamilies}; \texttt{AffineCoverJoin} and
\texttt{AffineRefinement} extend refinement to every compact open of a qcqs
scheme. These proofs include empty covers and introduce no new
foundational parameters. These formed the first checkpoint, checked with
114 modules and zero failures.

The comparison is now checked, including both inverse laws.
\texttt{MapLaws} supplies identity and inverse equations;
\texttt{Spaces} packages the existing local maps in geometric order,
without asserting univalence of the space type.
\texttt{AffineTests} proves the identity and composition laws of
$\Spec$ and the correspondence between affine maps and ring maps.
\texttt{FunctorPoints} defines the ordinary functor of points and its
natural transformations on all small rings.

\texttt{AffineEvaluation} proves evaluation for a spectrum, and
\texttt{AffineSource} transfers it to any affine source.
\texttt{OpenMaps}, \texttt{OpenAgreement}, and \texttt{OpenPullback}
provide the inclusion triangles and overlap equations used in the proof.
\texttt{TopOpen}, \texttt{AffineQcqs}, and \texttt{SliceQcqs}
give the qcqs properties needed for affine probes into compact opens.
\texttt{ProbeSeparation} proves that these probes detect equality.
\texttt{AtlasMap} constructs the recovered map from an actual affine
atlas, and \texttt{AtlasMapPoints} proves its equation on all affine tests.
\texttt{PointsUniversal.Over.equivalence} eliminates the merely given
atlas into the proposition of solutions and proves absolute full faithfulness.

For the relative result, \texttt{BasePoints} defines all $k$-algebra
tests, points over $k$, and their reindexing laws. It also proves the
equivalence between total relative points and ordinary points on a fixed
test ring. \texttt{RelativeFamilies} extends a relative natural family
to rings by canonical lifts, as in Section~5.
\texttt{RelativeUniversal.Over.equivalence} is the statement of
Theorem~\ref{thm:main}. Both morphism types and natural-transformation
types are proved to be sets. Morphisms lie in $\UU_\ell$ and the type of
natural transformations lies in $\UU_{\ell+1}$; the equivalence is stated
at the latter level. There is no resizing or finite-presentation
assumption on the tests.

The full tree checks with 133 modules and zero failures. The opening
audits are recorded in \texttt{descent.opens} and \texttt{points.opens}.
The remaining comparison problem is the essential-image theorem:
characterising functors admitting open affine atlases and reconstructing
schemes from that data. No such characterisation, or equivalence with
all set-valued functors, is asserted here.
\end{document}
